\subsection{Teorías del Aprendizaje Aplicadas a la Enseñanza de Compiladores}

El diseño pedagógico del prototipo se fundamenta en tres corrientes teóricas que la literatura especializada ha identificado como relevantes para la enseñanza de contenidos de alta abstracción en ciencias de la computación: la teoría de la carga cognitiva, la visualización de programas como estrategia de apoyo al aprendizaje, y la evaluación por pasos como mecanismo para hacer observable el proceso de interpretación.

\subsubsection{Teoría de la carga cognitiva}

La teoría de la carga cognitiva, desarrollada por Sweller, establece que la memoria de trabajo posee una capacidad limitada para procesar información nueva de forma simultánea \cite{sweller2024cognitive}. La teoría distingue tres tipos de carga: la intrínseca, determinada por la complejidad inherente del contenido y la interacción entre sus elementos; la extrínseca, generada por la forma en que la información se presenta; y la germana, asociada a los procesos mentales que contribuyen a la formación de esquemas de conocimiento duraderos. El diseño instruccional efectivo busca mantener la carga intrínseca en un nivel manejable, minimizar la carga extrínseca y maximizar la germana.

En el contexto de la enseñanza de compiladores e intérpretes, la carga intrínseca es elevada porque los conceptos involucrados, como gramáticas, AST, ambientes y clausuras, son abstractos y están interrelacionados de forma no trivial \cite{zagami2023}. Esta complejidad hace especialmente necesario controlar la carga extrínseca mediante decisiones de diseño que faciliten la interpretación de la información presentada al estudiante. Cuando el estudiante debe coordinar fuentes dispersas de información durante su trabajo, como una especificación de gramática en un archivo, el código del evaluador en otro y la salida del intérprete en la consola del sistema operativo, la presión sobre la memoria de trabajo aumenta y los recursos cognitivos disponibles para la comprensión del contenido disminuyen \cite{sweller2024cognitive}. El diseño de herramientas de apoyo para este tipo de cursos debe buscar reducir esa carga extrínseca integrando en un único entorno las representaciones relevantes del proceso de interpretación.

\subsubsection{Visualización de programas como apoyo al aprendizaje}

La visualización de programas es un área de investigación que estudia el uso de representaciones gráficas para apoyar la comprensión de conceptos abstractos de programación. Sorva et al. realizaron una revisión sistemática de los sistemas de visualización de programas para la enseñanza introductoria y concluyeron que la efectividad de estas herramientas está directamente relacionada con el grado en que la representación visual se corresponde con el modelo conceptual que el estudiante debe construir \cite{Sorva2013}. Cuando la visualización refleja con fidelidad las estructuras que el docente describe en clase, el estudiante puede establecer correspondencias directas entre la teoría y la herramienta, favoreciendo la comprensión y reduciendo el riesgo de formarse concepciones erróneas.

Python Tutor, desarrollado por Guo \cite{Guo2013}, es uno de los sistemas de visualización de ejecución más utilizados en educación en computación. La herramienta representa paso a paso el estado del ambiente de ejecución, incluyendo los marcos de activación, el montículo de objetos y los punteros entre ellos, para lenguajes como Python, Java y JavaScript. Su diseño demostró que la representación explícita del ambiente de ejecución facilita la comprensión de conceptos como el alcance de variables, la recursión y las referencias entre objetos. Sin embargo, Python Tutor no admite Racket ni visualiza la estructura interna de los intérpretes definidos por el usuario, lo que lo hace insuficiente para las necesidades específicas de la asignatura FLP.

La herramienta presentada por Spigariol et al. \cite{Spigariol2023Aprendiendo} constituye el antecedente más cercano al prototipo desarrollado en este trabajo en cuanto a objetivo pedagógico: combina un intérprete de un lenguaje funcional con componentes gráficos que visualizan las estructuras intermedias del proceso de interpretación. No obstante, dicha herramienta está diseñada para un lenguaje predefinido y no permite al estudiante definir su propia gramática ni trabajar con el modelo de EOPL y SLLGEN.

\subsubsection{Evaluación por pasos como estrategia pedagógica}

Clements, Flatt y Felleisen propusieron el modelo del \textit{stepper} algebraico como herramienta para hacer observable el proceso de reducción de expresiones en lenguajes funcionales \cite{Clements2001}. En este modelo, la evaluación de una expresión se presenta como una secuencia de pasos de reducción, donde cada paso muestra el estado del programa antes y después de la aplicación de una regla semántica. Este mecanismo permite al estudiante observar la correspondencia entre las reglas del evaluador y el comportamiento concreto del intérprete expresión por expresión.

La presentación de la evaluación como una secuencia de pasos discretos tiene además fundamento en la teoría de la carga cognitiva: al exponer un único paso de reducción a la vez, la herramienta controla la cantidad de información nueva que el estudiante debe procesar de forma simultánea, reduciendo la presión sobre la memoria de trabajo y favoreciendo la comprensión progresiva del comportamiento del intérprete \cite{sweller2024cognitive}. El prototipo desarrollado en este trabajo implementa este principio mediante el modo de evaluación paso a paso descrito en el Capítulo 4.

\subsection{Principios de Diseño de Herramientas Educativas Interactivas}

\subsubsection{Usabilidad en tecnología educativa}

La usabilidad se define como el grado en que un sistema puede ser utilizado por usuarios específicos para alcanzar objetivos concretos con efectividad, eficiencia y satisfacción en un contexto de uso determinado. En el ámbito de las herramientas educativas digitales, investigaciones recientes señalan que la usabilidad no incide únicamente en la eficiencia técnica sino también en la motivación y la retención del aprendizaje: una interfaz que genera fricción desplaza los recursos cognitivos del usuario desde la comprensión del contenido hacia la resolución de problemas de navegación o interacción \cite{lu2022usability}. Esta relación es especialmente relevante en herramientas para la enseñanza de compiladores, donde la carga cognitiva intrínseca ya es elevada.

En el ámbito de los cursos de compiladores, investigaciones previas han documentado que los estudiantes perciben las herramientas de desarrollo convencionales como una fuente adicional de dificultad, independientemente de la solidez técnica de los contenidos \cite{zagami2023}. Cuando la interfaz de una herramienta genera fricción, los recursos cognitivos se desvían desde la comprensión del contenido hacia la resolución de problemas de interacción, lo que reduce la capacidad disponible para el aprendizaje \cite{lu2022usability}. Por este motivo, el diseño de herramientas educativas para compiladores debe priorizar la minimización de los cambios de contexto y la integración de retroalimentación inmediata sobre los errores del estudiante.

\subsubsection{Organización visual y presentación simultánea de representaciones}

Norman señala que el diseño de interfaces debe respetar los límites de la atención humana, presentando la información de forma que la percepción y la interpretación requieran el mínimo esfuerzo posible \cite{norman2013design}. En el diseño de herramientas educativas para compiladores, este principio se traduce en la presentación simultánea de representaciones complementarias del proceso de interpretación, de modo que el estudiante pueda relacionar el código fuente con el AST y el ambiente sin necesidad de cambiar de ventana ni de contexto mental.

La literatura sobre herramientas educativas para compiladores confirma que la presentación simultánea de representaciones sincronizadas favorece la comprensión \cite{Stamenkovic2024AWE, jovanovic2021comvis}. Además, la diferenciación visual mediante color y estructura tipográfica permite al estudiante identificar la naturaleza de cada elemento de la interfaz sin necesidad de leer su contenido completo, lo que reduce el esfuerzo de interpretación y contribuye a mantener los recursos cognitivos disponibles para el análisis del comportamiento del intérprete \cite{Hannebauer2018}.
